% !TeX root = main.tex
% !TeX program = xelatex
% Translation of the current Chinese introduction.

\section{Introduction}
\label{sec:introduction}

Multi-agent communication requires authentication of both parties and control over communication initiation and subsequent requests. SAGA (Security Architecture for Governing Agentic Systems)~\cite{syros2025saga} is a security architecture for AI (artificial intelligence) agents. Its contact policies determine whether a sender may retrieve a receiver's OTK (one-time public key) to obtain an ACT (access control token). Subsequent requests are constrained by the ACT's identity binding, validity period, and request quota.

The original SAGA study used ProVerif~\cite{blanchet2016proverif} to analyze agent--Provider authentication, inter-agent authentication, and ACT confidentiality~\cite{syros2025saga}. These properties do not guarantee correct authorization enforcement. Rule selection, OTK issuance and use, and concurrent ACT quota updates require separate analysis.

Using SAGA's protocol specification and reference implementation~\cite{syros2025sagaV2,sagaReferenceImpl2026}, we construct separate Tamarin~\cite{meier2013tamarin} models for contact-policy matching and OTK issuance and use. The policy model captures rule selection and subsequent authorization to check rule priorities and consistency with expected decisions. An independent evaluator and policy-evaluation results from the reference implementation cross-check its outputs. The OTK model checks whether keys are used by the senders to which they were issued and whether issuance limits constrain successful use of distinct OTKs. A TLA$^{+}$ (Temporal Logic of Actions) model~\cite{lamport2002specifying} checks ACT request quotas under concurrency. Reference-implementation experiments assess the effects of incorrect policy matching and ACT quota violations. We identify three issues:

(1) The matcher fails to select the most specific matching rule, allowing a sender that should be denied access to obtain an ACT.

(2) OTKs are not bound to the senders to which they are issued. An authorized sender can therefore obtain new ACTs using OTKs issued to and forwarded by other senders after exhausting its own retrieval quota.

(3) ACT validation and quota deduction do not form a single atomic operation. Concurrent connections can therefore cause the receiver to process more requests than the token's initial quota permits.
